\section*{Chapitre \ref{ch:s1:momentum} --- Momentum}

\begin{solution}{exo:s1:momentum:1}
$0.12/0.40 = 0.3$~; une perte brute de 34~\% devient $0.3 \times (-0.34) = -10.2~\%$. C'est presque exactement avril 2009~: le poids était de 0.29 et la perte normalisée par la volatilité de 10.1~\%.
\end{solution}

\begin{solution}{exo:s1:momentum:2}
Un gain de $1/0.5 - 1 = 100~\%$. Après une perte de 50~\% et un gain de 100~\%, la stratégie est revenue à son point de départ~: $0.5 \times 2 - 1 = 0$. Des pertes de la taille de celles du momentum
demandent des années de rendement moyen pour s'en remettre, ce qui explique que son drawdown maximal (78.4~\%) compte davantage que sa moyenne.
\end{solution}

\begin{solution}{exo:s1:momentum:3}
Le poids échangé vaut deux fois le taux de rotation à sens unique~: $2 \times 5.7 \times 0.001 = 1.14~\%$ par an.
\end{solution}

\begin{solution}{exo:s1:momentum:4}
Pendant une baisse, le rendement de chaque action vaut à peu près son bêta multiplié par la chute du marché, plus son rendement propre~; les actions qui ont le plus chuté, les perdantes passées, ont
un bêta nettement plus élevé. Les vendre à découvert rend le portefeuille vendeur en bêta. Quand le marché rebondit fortement, les perdantes à fort bêta sont celles qui montent le plus, et le portefeuille perd
sur ses positions vendeuses plus qu'il ne gagne sur ses gagnantes à bêta plus faible, au moment où la volatilité est la plus élevée.
\end{solution}

\begin{solution}{exo:s1:momentum:5}
Sur un mois, les rendements s'inversent (chapitre~2). Inclure le dernier mois reviendrait à faire acheter les gagnantes du mois passé, qui tendent à en rendre une partie,
et diluerait le momentum par un pari de retournement de signe opposé.
\end{solution}

\begin{solution}{exo:s1:momentum:6}
La normalisation réduit l'exposition quand la volatilité récente est élevée, c'est-à-dire quand surviennent la plupart des krachs. Elle ne voit pas un krach venant d'un état calme~: en septembre
1939, la volatilité glissante du momentum était faible, le poids valait 0.92, et le krach est arrivé à presque pleine taille. La normalisation gère le risque prévisible, pas tout le risque.
\end{solution}

\begin{solution}{exo:s1:momentum:7}
Le ratio de Sharpe du portefeuille résiduel passe de 0.67 à 0.70 sur l'échantillon normalisé, celui du portefeuille en rendement total reste à 0.52. Sur le facteur français, il passe de 0.45
à 0.87. La normalisation aide d'autant plus que la volatilité de la stratégie varie et de manière prévisible~: celle du facteur français évolue d'un facteur
de plusieurs entre les années calmes et les crises, alors que celle des portefeuilles synthétiques varie beaucoup moins.
\end{solution}

\begin{solution}{exo:s1:momentum:8}
L'effet de levier multiplie aussi les pertes~: avec un levier de trois, la perte du facteur brut en 2009 aurait dépassé 100~\% du capital, et trois fois son drawdown maximal
de 78~\% équivaut à la ruine. Un ratio de Sharpe de 0.45 avec une asymétrie de $-3$ ne se démultiplie pas comme un rendement normal. Commencer par normaliser par la volatilité, puis choisir
le levier sur la série normalisée, et garder les années de krach dans l'échantillon.
\end{solution}

\begin{solution}{pb:s1:momentum:1}
\begin{enumerate}
\item Les actions ayant eu de forts rendements au cours des derniers mois continuent de surperformer~; le signal 12--1 est le rendement entre il y a douze mois et il y a un mois.
\item Rendements positifs et significatifs de l'achat des gagnantes et de la vente des perdantes sur des périodes de détention de 3 à 12 mois, non expliqués par le risque systématique, en partie
inversés dans les deux années suivantes.
\item 7.4~\% par an, volatilité de 16.3~\%, ratio de Sharpe de 0.45, asymétrie $-3.02$.
\item $-52.6~\%$ en août 1932~; $-78.4~\%$ de juin 1932 à septembre 1939.
\item Classement sur les résidus du modèle factoriel~; classement des secteurs sur leurs rendements. Le \omterm{def:s1:momentum:residual}{momentum résiduel} a rapporté environ deux fois plus de profits ajustés du risque~; le \omterm{def:s1:momentum:residual}{momentum sectoriel}
explique une grande partie du momentum des actions et était très rentable à lui seul.
\item Coefficients d'information (IC) de 0.020, 0.023 et 0.001~; ratios de Sharpe après coûts de 0.31, 0.45 et $-0.23$.
\item La même dérive avec moins de bruit factoriel~: une volatilité de 11.0~\% contre 13.4~\%.
\item Le marché synthétique ne contient aucun \omterm{def:s1:momentum:residual}{momentum sectoriel}~: ses facteurs sectoriels sont indépendants d'un jour à l'autre.
\item De courtes périodes de pertes très importantes, dans des états de panique après des baisses de marché et avec une forte volatilité, au moment où le marché rebondit.
\item Après une baisse, les perdantes sont des actions à fort bêta~; le portefeuille est vendeur en bêta, et un rebond le frappe.
\item Août et juillet 1932, avril 2009, septembre 1939, janvier 2001, juin 1938, juin 1931, avril 1933, novembre 2002 et janvier 2023.
\item Sur les deux années précédentes, le marché n'a baissé que pendant 6.7~\% des jours~: trop peu de marchés baissiers.
\item \textbf{Résultat nommé.} En 2009, le facteur momentum français a perdu 52.8~\% (49.4~\% de mars à mai)~; normalisé à 12~\% selon sa volatilité glissante sur 126 jours, il
a perdu 16.8~\% (16.5~\%).
\item Le ratio de Sharpe passe de 0.45 à 0.87 et l'asymétrie de $-3.02$ à $-0.40$.
\item Quand un krach naît d'un état calme~: septembre 1939, poids de 0.92, une perte normalisée de 29.1~\%.
\item Recourir à l'effet de levier chaque fois que la volatilité du momentum est inférieure à la cible, avec le financement et la capacité nécessaires.
\item Le rendement de la première demi-heure du marché prédit celui de la dernière demi-heure, dans les données de l'ETF S\&P 500 de 1993 à 2013 et dans dix autres ETF.
\item Le bêta du portefeuille et sa tendance après les baisses de marché, la volatilité réalisée de la stratégie, les mesures d'encombrement (crowding) et la rotation face aux coûts.
\item Peu~: environ 1.14~\% par an, contre 8--38~\% pour les portefeuilles quotidiens.
\item Des pertes rares et très importantes lorsque le marché rebondit après une baisse, que le rendement moyen ne laisse pas voir.
\end{enumerate}
\end{solution}

\begin{solution}{iq:s1:momentum:1}
Acheter les actions qui ont surperformé au cours des douze derniers mois et vendre celles qui ont sous-performé. Le dernier mois est écarté, car à un horizon
d'un mois, les rendements s'inversent.
\end{solution}

\begin{solution}{iq:s1:momentum:2}
Après les baisses de marché, les perdantes que le portefeuille détient à découvert sont des actions à fort bêta~; quand le marché rebondit fortement, souvent dans des états de panique à forte volatilité, elles montent
le plus et le portefeuille perd lourdement~: 1932, 2009.
\end{solution}

\begin{solution}{iq:s1:momentum:3}
La normaliser par sa volatilité prévue~; couvrir ou encadrer par un cap son bêta de marché~; la réduire après les baisses de marché à forte volatilité (état de panique de Daniel et Moskowitz)~;
utiliser le \omterm{def:s1:momentum:residual}{momentum résiduel}~; diversifier avec la value et d'autres stratégies qui réussissent lors des rebonds.
\end{solution}

\begin{solution}{iq:s1:momentum:4}
Le momentum en rendement total porte les gagnants factoriels de l'année précédente, si bien que ses expositions aux facteurs fluctuent~; le \omterm{def:s1:momentum:residual}{momentum résiduel} classe les actions sur leur part spécifique et se rapproche
de la neutralité factorielle. Le \omterm{def:s1:momentum:residual}{momentum résiduel}, pour un risque plus stable et, d'après les résultats publiés et ceux obtenus ici, pour un rendement ajusté du risque plus élevé.
\end{solution}

\begin{solution}{iq:s1:momentum:5}
Le marché a baissé et les perdantes que le portefeuille détient à découvert sont désormais des actions à fort bêta~: le portefeuille est devenu vendeur du marché. Couvrir le bêta ou réduire le portefeuille,
et surveiller un rebond, moment où surviennent les krachs.
\end{solution}

\begin{solution}{iq:s1:momentum:6}
Avec le rendement $r_t = \mu + \sigma_t\epsilon_t$ et le poids $w_t = c/\sigma_t$, le rendement normalisé a pour moyenne $c\mu\,E[1/\sigma_t]$ et pour variance $c^2$ (auquel s'ajoute un terme
issu de $\mu^2$), donc son ratio de Sharpe vaut environ $\mu E[1/\sigma_t]$, qui, par l'inégalité de Jensen, dépasse $\mu/\sqrt{E[\sigma_t^2]}$, le ratio brut. La normalisation abaisse le
ratio lorsque la moyenne est élevée au moment précis où la volatilité l'est aussi (la normalisation réduit alors l'exposition quand la stratégie rapporte le plus), ou lorsque la prévision de la volatilité est médiocre.
\end{solution}
